% =====================================================================
\chapter{Extremos libres y condicionados}
\label{tema:5}
% =====================================================================

Entrenar un modelo es minimizar una función de muchas variables, y estimar probabilidades por máxima verosimilitud es maximizar una función con la restricción de que las probabilidades sumen uno. Este tema estudia ambos problemas en dimensión dos y tres: los extremos libres, que se localizan donde se anula el gradiente y se clasifican con la matriz hessiana, y los extremos condicionados, que se resuelven con multiplicadores de Lagrange. La presentación sigue las sesiones~28--31 y~39--40 de 18.02SC~\cite{mit1802}, las §4.7--4.8 de OpenStax~\cite{openstax3}, las §9.9--9.12 y~9.14 de Apostol~\cite{apostol2}, el texto de Marsden y Tromba~\cite{marsden} y las §13.6--13.7 de Strang~\cite{strangcalc}.

\section{Extremos locales y puntos críticos}

\begin{definicion}[Extremos]\label{def:extremos}
Sea $f:D\subseteq\R^n\to\R$ y $\va\in D$. Se dice que $f$ tiene en $\va$ un \emph{máximo local} si existe $\delta>0$ tal que $f(\vx)\le f(\va)$ para todo $\vx\in D$ con $\norm{\vx-\va}<\delta$; un \emph{mínimo local}, con $\ge$. Si la desigualdad vale para todo $\vx\in D$, el extremo es \emph{absoluto} (o global). Un punto $\va$ interior al dominio es \emph{crítico} si $\nabla f(\va)=\vcero$ o si alguna derivada parcial no existe en $\va$.
\end{definicion}

\begin{teorema}[Condición necesaria de Fermat]\label{teo:fermat}
Si $f$ tiene un extremo local en un punto interior $\va$ y sus derivadas parciales existen en $\va$, entonces $\nabla f(\va)=\vcero$.
\end{teorema}

\begin{proof}
Sea $i\in\{1,\dots,n\}$ y $\varphi(t)=f(\va+t\mathbf e_i)$, definida para $t$ pequeño porque $\va$ es interior. La función de una variable $\varphi$ tiene un extremo local en $t=0$ y es derivable allí con $\varphi'(0)=f_{x_i}(\va)$. Por el teorema de Fermat de una variable, $\varphi'(0)=0$.
\end{proof}

El recíproco es falso: $f(x,y)=x^2-y^2$ tiene $\nabla f(0,0)=\vcero$ y no tiene extremo en el origen, pues crece a lo largo del eje $x$ y decrece a lo largo del eje $y$. Un punto crítico que no es extremo local se llama \emph{punto de silla}.

\section{Clasificación: el criterio del hessiano}

\begin{definicion}[Matriz hessiana]\label{def:hessiana}
Si $f(x,y)$ tiene derivadas segundas, su \emph{matriz hessiana} y su \emph{discriminante} en $\va$ son
\[
H_f(\va)=\begin{pmatrix}f_{xx}&f_{xy}\\f_{yx}&f_{yy}\end{pmatrix}(\va),\qquad
\Delta(\va)=\det H_f(\va)=f_{xx}f_{yy}-f_{xy}^2 ,
\]
donde se ha usado el teorema de Schwarz para escribir $f_{xy}=f_{yx}$.
\end{definicion}

La idea del criterio es la fórmula de Taylor de segundo orden: si $f$ tiene derivadas segundas continuas y $\nabla f(\va)=\vcero$, entonces, con $\mathbf h=(h,k)$,
\begin{equation}\label{eq:t5-taylor}
f(\va+\mathbf h)=f(\va)+\tfrac12\,Q(\mathbf h)+\varepsilon(\mathbf h)\norm{\mathbf h}^2,\qquad Q(h,k)=Ah^2+2Bhk+Ck^2,
\end{equation}
con $A=f_{xx}(\va)$, $B=f_{xy}(\va)$, $C=f_{yy}(\va)$ y $\varepsilon(\mathbf h)\to0$. La demostración de~\eqref{eq:t5-taylor} aplica la fórmula de Taylor de una variable a $t\mapsto f(\va+t\mathbf h)$ (Apostol~\cite[§9.9--9.12]{apostol2}; módulo~4, tema~4). Cerca de un punto crítico, por tanto, el signo de $f(\va+\mathbf h)-f(\va)$ lo decide la \emph{forma cuadrática} $Q$.

\begin{lema}[Signo de una forma cuadrática]\label{lem:forma}
Sea $Q(h,k)=Ah^2+2Bhk+Ck^2$ y $\Delta=AC-B^2$.
\begin{enumerate}
  \item Si $\Delta>0$ y $A>0$, entonces $Q(h,k)>0$ para todo $(h,k)\neq(0,0)$.
  \item Si $\Delta>0$ y $A<0$, entonces $Q(h,k)<0$ para todo $(h,k)\neq(0,0)$.
  \item Si $\Delta<0$, $Q$ toma valores positivos y negativos.
\end{enumerate}
\end{lema}

\begin{proof}
Si $A\neq0$, completando el cuadrado,
\[
Q(h,k)=A\Bigl(h+\frac BAk\Bigr)^2+\frac{AC-B^2}{A}k^2=\frac1A\Bigl[(Ah+Bk)^2+\Delta k^2\Bigr].
\]
(i) Si $\Delta>0$ y $A>0$, el corchete es suma de dos términos no negativos que solo se anulan a la vez si $k=0$ y $Ah+Bk=Ah=0$, es decir, si $(h,k)=(0,0)$. (ii) Es (i) aplicado a $-Q$. (iii) Si $\Delta<0$ y $A\neq0$: con $k=0$, $h=1$, $Q=A$; con $k=1$, $h=-B/A$, $Q=\Delta/A$, de signo contrario a $A$. Si $A=0$, entonces $\Delta=-B^2<0$ obliga a $B\neq0$, y $Q(h,1)=2Bh+C$ toma ambos signos al variar $h$.
\end{proof}

\begin{teorema}[Criterio del hessiano]\label{teo:hessiano}
Sea $f$ con derivadas segundas continuas cerca de un punto crítico $\va$ con $\nabla f(\va)=\vcero$, y sea $\Delta=\Delta(\va)$.
\begin{enumerate}
  \item Si $\Delta>0$ y $f_{xx}(\va)>0$, $f$ tiene un mínimo local en $\va$.
  \item Si $\Delta>0$ y $f_{xx}(\va)<0$, $f$ tiene un máximo local en $\va$.
  \item Si $\Delta<0$, $\va$ es un punto de silla.
  \item Si $\Delta=0$, el criterio no decide.
\end{enumerate}
\end{teorema}

\begin{proof}[Idea de la demostración]
(i) Por el \cref{lem:forma}, $Q>0$ en la circunferencia unidad, que es cerrada y acotada; por el teorema de Weierstrass (\cref{teo:weierstrass}) alcanza allí un mínimo $m>0$, y por homogeneidad $Q(\mathbf h)\ge m\norm{\mathbf h}^2$. En~\eqref{eq:t5-taylor}, si $\norm{\mathbf h}$ es tan pequeño que $\abs{\varepsilon(\mathbf h)}<m/4$, se obtiene $f(\va+\mathbf h)-f(\va)\ge\frac m4\norm{\mathbf h}^2>0$. (ii) se reduce a (i) con $-f$. (iii) Si $Q(\mathbf h_1)>0>Q(\mathbf h_2)$, a lo largo de las rectas $t\mathbf h_1$ y $t\mathbf h_2$ la función crece y decrece respectivamente para $t$ pequeño, por el mismo argumento. (iv) Lo muestran los ejemplos del \cref{eje:5-degenerado}.
\end{proof}

\begin{ejemplo}\label{ej:t5-hessiano}
(a) $f(x,y)=x^2+xy+y^2-3x$: $f_x=2x+y-3$ y $f_y=x+2y$ se anulan en $(2,-1)$. La hessiana es constante, $\begin{pmatrix}2&1\\1&2\end{pmatrix}$, con $\Delta=3>0$ y $f_{xx}=2>0$: mínimo local, de valor $f(2,-1)=-3$.

(b) $g(x,y)=x^3-3x+y^2$: $g_x=3x^2-3$ y $g_y=2y$ se anulan en $(\pm1,0)$. Como $H_g=\begin{pmatrix}6x&0\\0&2\end{pmatrix}$, en $(1,0)$ es $\Delta=12>0$ con $g_{xx}=6>0$ (mínimo local, $g=-2$), y en $(-1,0)$ es $\Delta=-12<0$ (punto de silla).
\end{ejemplo}

\section{Extremos absolutos en regiones cerradas y acotadas}

\begin{teorema}[Weierstrass]\label{teo:weierstrass}
Una función continua sobre un conjunto cerrado y acotado de $\R^n$ alcanza en él un máximo y un mínimo absolutos.
\end{teorema}

La demostración usa la completitud de los números reales y se omite; para una variable es el teorema de los valores extremos del cálculo elemental. En la práctica, para hallar los extremos absolutos de $f$ en una región $R$ cerrada y acotada: (1) se hallan los puntos críticos interiores; (2) se estudia $f$ sobre la frontera, que suele reducirse a funciones de una variable o a un problema con multiplicadores de Lagrange; (3) se comparan todos los valores obtenidos.

\begin{ejemplo}\label{ej:t5-absolutos}
Extremos de $f(x,y)=x^2+y^2-2x$ en el disco $x^2+y^2\le4$. Interior: $\nabla f=(2x-2,2y)=\vcero$ en $(1,0)$, con $f(1,0)=-1$. Frontera: si $x^2+y^2=4$, $f=4-2x$ con $x\in[-2,2]$, que varía entre $0$ (en $(2,0)$) y $8$ (en $(-2,0)$). Comparando: el máximo absoluto es $8$, en $(-2,0)$, y el mínimo absoluto es $-1$, en $(1,0)$.
\end{ejemplo}

\section{Multiplicadores de Lagrange}

Se busca el máximo o el mínimo de $f(x,y)$ cuando $(x,y)$ está obligado a moverse sobre la curva $g(x,y)=c$. En un extremo, la curva de nivel de $f$ que pasa por el punto no puede cruzar la restricción, porque entonces se podría seguir avanzando sobre $g=c$ hacia valores mayores y menores de $f$: ambas curvas deben ser \emph{tangentes}, y sus gradientes, paralelos (\cref{fig:t5-lagrange}).

\begin{figure}[ht]
\centering
\begin{tikzpicture}[scale=1.15,>={Stealth[length=2mm]},font=\small]
  \draw[->] (-1.9,0) -- (2.4,0) node[right] {$x$};
  \draw[->] (0,-1.9) -- (0,2.4) node[above] {$y$};
  \draw[thick] (0,0) circle ({sqrt(2)});
  \foreach \c in {0,1,2,3} \draw[gray,dashed] ({\c-2.2},2.2) -- (2.2,{\c-2.2});
  \fill (1,1) circle (1.6pt);
  \draw[->,thick] (1,1) -- (1.55,1.55) node[above right] {$\nabla f\parallel\nabla g$};
  \node at (-1.25,-1.45) {$g=2$};
  \node[gray] at (2.15,-0.45) {$f=c$};
\end{tikzpicture}
\caption{Máximo de $f(x,y)=x+y$ sobre la circunferencia $g(x,y)=x^2+y^2=2$. Las rectas de nivel de $f$ (discontinuas) cortan a la circunferencia salvo la de nivel $2$, que es tangente en $(1,1)$; allí $\nabla f=(1,1)$ y $\nabla g=(2,2)$ son paralelos.}
\label{fig:t5-lagrange}
\end{figure}

\begin{teorema}[Multiplicadores de Lagrange]\label{teo:lagrange}
Sean $f$ y $g$ funciones de $(x,y)$ con derivadas parciales continuas, y sea $\va$ un punto de la curva $g=c$ con $\nabla g(\va)\neq\vcero$. Si $f$, restringida a la curva, tiene un extremo local en $\va$, existe $\lambda\in\R$ (el \emph{multiplicador}) tal que
\[
\nabla f(\va)=\lambda\nabla g(\va).
\]
El mismo enunciado vale para $f$ y $g$ de tres variables y la superficie $g=c$.
\end{teorema}

\begin{proof}
Como $\nabla g(\va)\neq\vcero$, el teorema de la función implícita~\cite[§9.6--9.7]{apostol2} asegura que, cerca de $\va$, la curva $g=c$ es la trayectoria de una curva derivable $\vr(t)$ con $\vr(0)=\va$ y $\vr\,'(0)\neq\vcero$. La función de una variable $t\mapsto f(\vr(t))$ tiene un extremo local en $t=0$, luego, por la regla de la cadena, $0=\nabla f(\va)\cdot\vr\,'(0)$. Además $\nabla g(\va)\cdot\vr\,'(0)=0$ por el \cref{teo:gradiente-normal}. En el plano, los vectores perpendiculares al vector no nulo $\vr\,'(0)$ forman una recta, de modo que $\nabla f(\va)$ y $\nabla g(\va)$ son proporcionales; como $\nabla g(\va)\neq\vcero$, $\nabla f(\va)=\lambda\nabla g(\va)$. En $\R^3$ el argumento se aplica a todas las curvas de la superficie que pasan por $\va$: $\nabla f(\va)$ es perpendicular a su plano tangente y, por tanto, paralelo a la normal $\nabla g(\va)$.
\end{proof}

El método consiste en resolver el sistema formado por $\nabla f=\lambda\nabla g$ y $g=c$ (tres ecuaciones con incógnitas $x,y,\lambda$ en el plano; cuatro en el espacio). Sus soluciones son los \emph{candidatos}; si la restricción es cerrada y acotada, el teorema de Weierstrass garantiza que el mayor valor de $f$ entre ellos es el máximo y el menor, el mínimo.

\begin{ejemplo}\label{ej:t5-lagrange-circulo}
Extremos de $f=x+y$ sobre $x^2+y^2=2$. Las ecuaciones $1=2\lambda x$, $1=2\lambda y$ dan $x=y=1/(2\lambda)$, y la restricción, $2x^2=2$, luego $x=\pm1$. Candidatos: $(1,1)$, con $f=2$, y $(-1,-1)$, con $f=-2$. La circunferencia es cerrada y acotada: el máximo es $2$ y el mínimo $-2$.
\end{ejemplo}

\begin{ejemplo}\label{ej:t5-lagrange-plano}
El punto del plano $x+2y+2z=9$ más cercano al origen minimiza $f=x^2+y^2+z^2$ con $g=x+2y+2z=9$. Las ecuaciones $2x=\lambda$, $2y=2\lambda$, $2z=2\lambda$ dan $(x,y,z)=(\lambda/2,\lambda,\lambda)$; sustituyendo, $\tfrac{9}{2}\lambda=9$, $\lambda=2$ y el punto es $(1,2,2)$, a distancia~$3$ del origen, de acuerdo con la fórmula $\abs d/\norm{\vn}=9/3$ del \cref{eje:2-distancia}. Aquí la restricción no es acotada, pero $f$ crece sin límite al alejarse, de modo que el único candidato es el mínimo.
\end{ejemplo}

\begin{observacion}[Errores frecuentes]
\begin{enumerate}[label=(\alph*)]
\item $\nabla f=\vcero$ no implica extremo (puntos de silla); y con $\Delta=0$ no se puede concluir nada sin un análisis adicional.
\item El criterio del hessiano sirve para puntos críticos interiores; en la frontera de una región los extremos se buscan aparte.
\item Lagrange proporciona \emph{candidatos}, no extremos: hay que comparar sus valores y justificar la existencia del máximo o del mínimo.
\item Al resolver $\nabla f=\lambda\nabla g$ no se debe dividir por una variable que pueda anularse; se pierden soluciones (\cref{eje:5-hiperbola}).
\end{enumerate}
\end{observacion}

\begin{observacion}[Conexión con el curso]
El estimador de máxima verosimilitud de un modelo de bigramas (módulo~1, tema~3) maximiza $\sum_i n_i\ln p_i$ con la restricción $\sum_ip_i=1$; con un multiplicador de Lagrange resulta la frecuencia relativa $p_i=n_i/N$ (\cref{eje:5-mle}). El mismo método demuestra que la distribución de máxima entropía sobre $K$ valores es la uniforme (módulo~6). La hessiana y su relación con la convexidad y con la velocidad del descenso de gradiente se desarrollan en los módulos~4 (tema~4) y~7, y la recta de mínimos cuadrados (\cref{eje:5-minimos-cuadrados}) es el punto de partida de la regresión lineal del módulo~8.
\end{observacion}

\begin{lecturas}
\item MIT 18.02SC~\cite{mit1802}, unidad~2, parte~A, sesiones~28--31 (optimización, mínimos cuadrados, criterio de la derivada segunda), y parte~C, sesiones~39--40 (multiplicadores de Lagrange y su demostración).
\item OpenStax \emph{Calculus}, vol.~3~\cite{openstax3}, §4.7 «Maxima/Minima Problems» y §4.8 «Lagrange Multipliers».
\item Apostol~\cite{apostol2}, §9.9--9.12 (hessiana, criterio de la derivada segunda, fórmula de Taylor) y §9.14 (extremos condicionados y multiplicadores de Lagrange).
\item Strang, \emph{Calculus}~\cite{strangcalc}, §13.6 «Maxima, Minima, and Saddle Points» y §13.7 «Constraints and Lagrange Multipliers».
\end{lecturas}

\section*{Ejercicios del tema 5}
\addcontentsline{toc}{section}{Ejercicios del tema 5}

\begin{ejercicio}{1}{eje:5-conceptos}
Razone si son verdaderas o falsas: (a) si $\nabla f(\va)=\vcero$, entonces $f$ tiene un extremo local en $\va$; (b) si en un punto crítico $\Delta=0$, el punto es de silla; (c) toda solución del sistema de Lagrange es un máximo o un mínimo condicionado; (d) una función continua en el disco abierto $x^2+y^2<1$ alcanza siempre su máximo.
\end{ejercicio}

\begin{ejercicio}{2}{eje:5-clasificar-1}
Halle y clasifique los puntos críticos de: (a) $x^2-xy-2y^2-3x-3y+1$; (b) $3x^2+xy+y^2-x-2y+4$.
\fuente{adaptado de MIT 18.02SC~\cite{mit1802}, problemas suplementarios, 2H-1.}
\end{ejercicio}

\begin{ejercicio}{2}{eje:5-clasificar-2}
Halle y clasifique los puntos críticos de $f(x,y)=x^3-3xy+y^3$.
\fuente{adaptado de MIT 18.02SC~\cite{mit1802}, problemas suplementarios, 2H-1.}
\end{ejercicio}

\begin{ejercicio}{2}{eje:5-clasificar-3}
Halle y clasifique los puntos críticos de $f(x,y)=2x^4+y^2-xy+1$.
\fuente{adaptado de MIT 18.02SC~\cite{mit1802}, problemas suplementarios, 2H-1.}
\end{ejercicio}

\begin{ejercicio}{2}{eje:5-openstax}
Halle y clasifique los puntos críticos de: (a) $f(x,y)=x^3+2xy-2x-4y$; (b) $g(x,y)=4x^2+9y^2+8x-36y+24$. En (b), confirme el resultado completando cuadrados.
\fuente{adaptado de OpenStax \emph{Calculus}, vol.~3~\cite{openstax3}, §4.7, punto de control~4.34 y ejemplo~4.39.}
\end{ejercicio}

\begin{ejercicio}{2}{eje:5-absolutos}
Halle el máximo y el mínimo absolutos de $f(x,y)=xy-x-y+2$ en el cuadrado $0\le x\le2$, $0\le y\le2$. (Indicación: $f=(x-1)(y-1)+1$.)
\fuente{adaptado de MIT 18.02SC~\cite{mit1802}, problemas suplementarios, 2H-4.}
\end{ejercicio}

\begin{ejercicio}{2}{eje:5-disco}
Halle el máximo y el mínimo de $f(x,y)=x+\sqrt3\,y+2$ en el disco $x^2+y^2\le1$.
\fuente{adaptado de MIT 18.02SC~\cite{mit1802}, problemas suplementarios, 2H-5.}
\end{ejercicio}

\begin{ejercicio}{2}{eje:5-caja}
Una caja rectangular del primer octante tiene un vértice en el origen, caras sobre los planos coordenados y el vértice opuesto $P=(x,y,z)$ sobre el plano $x+2y+3z=18$. Halle con multiplicadores de Lagrange el $P$ que maximiza el volumen y justifique que es un máximo.
\fuente{adaptado de MIT 18.02SC~\cite{mit1802}, problemas suplementarios, 2I-1.}
\end{ejercicio}

\begin{ejercicio}{2}{eje:5-hiperbola}
Halle los extremos de $f(x,y)=x^2-y^2$ sobre la circunferencia $x^2+y^2=3$ con multiplicadores de Lagrange. Indique, para cada solución, el valor de $\lambda$ y si $\nabla f$ y $\nabla g$ tienen el mismo sentido o sentidos opuestos.
\fuente{adaptado de MIT 18.02SC~\cite{mit1802}, \emph{problem set}~6, parte~II, problema~1.}
\end{ejercicio}

\begin{ejercicio}{2}{eje:5-resistencias}
Una corriente $I>0$ se reparte entre dos resistencias en paralelo, $R_1$ y $R_2$, en corrientes $I_1+I_2=I$. La potencia disipada es $I_1^2R_1+I_2^2R_2$. Halle el reparto que la minimiza y compruebe que es el de las leyes de Kirchhoff ($I_1R_1=I_2R_2$).
\fuente{adaptado de MIT 18.02SC~\cite{mit1802}, \emph{problem set}~6, parte~II, problema~2.}
\end{ejercicio}

\begin{ejercicio}{2}{eje:5-minimos-cuadrados}
Halle la recta $y=ax+b$ que minimiza $D(a,b)=\sum_{i=1}^{3}\bigl(y_i-(ax_i+b)\bigr)^2$ para los puntos $(1,1)$, $(2,3)$, $(3,2)$, y compruebe con el criterio del hessiano que el punto crítico es un mínimo.
\fuente{adaptado de MIT 18.02SC~\cite{mit1802}, problemas suplementarios, 2G-1.}
\end{ejercicio}

\begin{ejercicio}{3}{eje:5-degenerado}
Demuestre que $f_1=x^4+y^4$, $f_2=-x^4-y^4$ y $f_3=x^4-y^4$ tienen todas un punto crítico en el origen con $\Delta=0$, y que en él $f_1$ tiene un mínimo, $f_2$ un máximo y $f_3$ un punto de silla. ¿Qué concluye sobre el caso~(iv) del \cref{teo:hessiano}?
\end{ejercicio}

\begin{ejercicio}{3}{eje:5-mle}
Sean $n_1,\dots,n_K>0$ los recuentos de $K$ símbolos en un corpus y $N=\sum_in_i$. Se quiere maximizar la log-verosimilitud $\ell(p_1,\dots,p_K)=\sum_{i=1}^{K}n_i\ln p_i$ con $p_i>0$ y $\sum_ip_i=1$. (a) Demuestre con un multiplicador de Lagrange que el único candidato es $\hat p_i=n_i/N$. (b) Demuestre que es el máximo absoluto usando la desigualdad $\ln t\le t-1$ para $t>0$.
\end{ejercicio}

\begin{ejercicio}{3}{eje:5-media-geometrica}
(a) Demuestre que el máximo de $xyz$ sobre $x+y+z=S$, $x,y,z\ge0$ ($S>0$ fijo) existe, se alcanza en el interior y vale $(S/3)^3$. (b) Deduzca la desigualdad entre las medias aritmética y geométrica: $\sqrt[3]{abc}\le\frac{a+b+c}{3}$ para $a,b,c\ge0$.
\end{ejercicio}

\begin{ejercicio}{3}{eje:5-cs-lagrange}
Sea $\va\in\R^3$ no nulo. (a) Halle con multiplicadores de Lagrange el máximo y el mínimo de $f(\vx)=\va\cdot\vx$ sobre la esfera unidad $\norm{\vx}=1$. (b) Deduzca de nuevo la desigualdad de Cauchy--Schwarz. (c) Interprete el resultado a la luz del \cref{cor:maximo-crecimiento}.
\end{ejercicio}
